
%\documentclass[./main.tex(path of main file)]{subfiles}
\documentclass[10pt]{article}
\usepackage{amsmath}
\DeclareMathOperator*{\argmin}{arg\,min} % thin space, limits underneath in displays
\DeclareMathOperator*{\argmax}{arg\,max} % thin space, limits underneath in displays
\newtheorem{thm}{Theorem}
\usepackage{amssymb}
\usepackage{amsfonts}
\usepackage{mathrsfs}
\usepackage{bm}
\usepackage{indentfirst}
\setlength{\parindent}{0em}
\usepackage[margin=1in]{geometry}
\usepackage{graphicx}
\usepackage{setspace}
\doublespacing
\usepackage[flushleft]{threeparttable}
\usepackage{booktabs,caption}
\usepackage{float}
\usepackage[sort,comma]{natbib}
\usepackage[hidelinks]{hyperref}
\usepackage{booktabs}
\usepackage{multirow}

\usepackage{import}
\usepackage{xifthen}
\usepackage{pdfpages}
\usepackage{transparent}
\usepackage{subfiles}
\usepackage{blindtext}
\usepackage{glossaries}
\usepackage{pdflscape}
\usepackage{adjustbox}
\usepackage{rotating}
\usepackage{eurosym,geometry,ulem,subcaption,caption,color,sectsty,comment,footmisc,pdflscape,array}

% Format definition list
\newenvironment{mylist}[1]%
  {\begin{list}{}{\settowidth{\labelwidth}{\textbf{#1}}
   \setlength{\leftmargin}{\labelwidth}
   \addtolength{\leftmargin}{\labelsep}
   \renewcommand{\makelabel}[1]{\textbf{\hfill##1}}}}%
  {\end{list}}



\newcommand{\incfig}[1]{%
\def\svgwidth{\columnwidth}
\import{./figures/}{#1.pdf_tex}
}




\title{}
\author{}
\date{}


\begin{document}
%\maketitle
\tableofcontents
\newpage

\section{Introduction}
{\textbf {Types of price rigidity}}:\\

Many prices are reviewed on a predetermined schedule and are only rarely changed at other times.
Thus price changes are not purely {\underline {state-dependent}} (that is, triggered by developments 
within the economy, regardless of the time over which the developments have occurred); they are partly
{\underline {time-dependent}} (that is, triggered by the passage of time).
\\


Section \ref{sec:DNK} presents a common framework for all the models of this part of the chapter.
Section \ref{sec:Fischer} through \ref{sec:Calvo} then consider three baseline models of 
{\underline {time-dependent}} price adjustment:
the Fischer, or Fischer-Phelps-Taylor model (Fischer, 1977; Phelps and Taylor, 1977); the Taylor
model (Taylor, 1979); and the Calvo model (Calvo, 1983). All three models posit that prices
(or wages) are set by multiperiod contracts or commitments. The central result of the models is that
multiperiod contracts lead to gradual adjustment of the price level to nominal disturbances. As a
result, {\underline {aggregate demand disturbances have persistent real effects.}}

The Taylor and Calvo models differ from the Fischer model in one important respect. The Fischer model
assumes that {\underline {prices are predetermined but not fixed}}. That is, when a multiperiod
contract sets prices for several periods, it can specify a different for each period. In the 
Taylor and Calvo models, in contrast, prices are fixed: a contract must specify the same price each
period it is in effect.

The difference between the Taylor and Calvo models is smaller. In the Taylor model, opportunities to
change prices arrive deterministically, and each price is in effect for the same number of periods.
In the Calvo model, opportunities to change prices arrive randomly, and so the number of periods
a price is in effect is stochastic.

Section \ref{sec:state-dependent} then turns to two baseline models of {\underline {state-dependent}}
price adjustment, the Caplin-Spulber and Danziger-Golosov-Lucas models (Caplin and Spulber, 1987;
Danziger, 1999; Golosov and Lucas, 2007) In both, the only barrier to price adjustment is a constant
fixed cost. There are two differences between the models. First, money growth is always positive in the Caplin-Spulber model, while the version of the Danziger-GolosovLucas model we will consider assumes no trend money growth. Second, the Caplin-Spulber model assumes no firm-specific shocks, while the DanzigerGolosov-Lucas model includes them. Both models deliver strong results about the effects of monetary disturbances, but for very different reasons.

Section \ref{sec:dynamic_price_adjustment} considers two more models of dynamic price adjustment:
the Calvo-with-indexation modoel and the Mankiw-Reis model (Mankiw and Reis, 2002;
Christiano, Eichenbaum, and Evans, 2005). These models are more complicated than the models of 
earlier sections, but appear to have more hope of fitting key facts about inflation dynamics.


\section{Building Blocks of Dynamic New Keynesian Models} \label{sec:DNK}
\section{The Fischer (or Fischer-Phelps-Taylor) Model} \label{sec:Fischer}
\section{The Taylor Model} \label{sec:Taylor}
\section{The Calvo Model} \label{sec:Calvo}
\section{State-Dependent Pricing} \label{sec:state-dependent}
\section{Dynamic Price Adjustment} \label{sec:dynamic_price_adjustment}














\bibliographystyle{plainnat}
\bibliography{my_bib}

\end{document}

